\section*{Chapter \ref{ch:s2:the-volatility-index-complex} --- The Volatility-Index Complex}

\begin{solution}{exo:s2:the-volatility-index-complex:1}
$A\,L(L - 1)\,r = 100 \times (-1) \times (-2) \times 0.3 = 60$: it buys futures worth 60.
\end{solution}

\begin{solution}{exo:s2:the-volatility-index-complex:2}
Twice the index: $100 \times 2 \times 1 \times 0.3 = 60$, also a purchase. Once the index: $L(L - 1) = 0$, nothing.
\end{solution}

\begin{solution}{exo:s2:the-volatility-index-complex:3}
$10w + 38(1 - w) = 30$ gives $w = 8/28 = 0.286$ on the first contract.
\end{solution}

\begin{solution}{exo:s2:the-volatility-index-complex:4}
The curve was in contango on 89\% of the year's days, so the long position paid the roll-down almost every day; VIX's rise came in spikes that
faded, and the index ended higher while the futures the long held had each converged down to it.
\end{solution}

\begin{solution}{exo:s2:the-volatility-index-complex:5}
The future prices the expected square root of a variance, and the square root is concave: by Jensen's inequality, the expected root is below the root
of the expected variance. The gap grows with the volatility of volatility.
\end{solution}

\begin{solution}{exo:s2:the-volatility-index-complex:6}
Their flow $A\,L(L - 1)\,r$ has the sign of the day's move for $L = -1$ and $L = 2$: after a rise in the futures they must buy more at the close,
which pushes the futures up further on the day when liquidity is thinnest.
\end{solution}

\begin{solution}{exo:s2:the-volatility-index-complex:7}
The index spends most of its time below its mean (median 17.7 against a mean of 19.6), so futures expecting reversion stand above it: contango on 65\%
of days with no premium at all. The long index's average daily return is close to zero ($-3.7\%$ a year), but its log return is $-26\%$ a year:
a series with 88\% volatility compounds far below its average.
\end{solution}

\begin{solution}{exo:s2:the-volatility-index-complex:8}
Returns do not scale with size in a short-volatility book: at half the capital the book was wiped out by the first crash, and at the full capital
sooner. The first crash's index rise of 189\% is more than the whole capital of a book short at 0.53 of it or more.
\end{solution}

\begin{solution}{pb:s2:the-volatility-index-complex:1}
\begin{enumerate}
\item The market's expectation of 30-day S\&P 500 volatility, from out-of-the-money SPX and SPXW options at two terms interpolated to 30 days.
\item The later future against the earlier or the index; contango on 85.0\% of days, the second contract 0.89 points above the first and the first
0.63 above VIX.
\item Each future converging down to the index as expiry nears; the long lost 99.99\% from 2013 to 2026, a log rate of 66.4\% a year.
\item 2014: $-25\%$ with VIX up 35\%; 2015: $-37\%$ with VIX up 2\%; 2025: $-45\%$ with VIX down 17\%.
\item A fund or note delivering a multiple of the daily return of a VIX futures index, rebalanced at the close.
\item Value $A(1 + Lr)$, futures held $LA(1 + r)$, target $LA(1 + Lr)$; the difference is $A\,L(L - 1)\,r$.
\item Inverse ($L = -1$) and leveraged ($L = 2$) products.
\item The most popular VIX products are virtually guaranteed to lose over time; those on the short-term indices had lost nearly \$4 billion.
\item VIX 17.31 to 37.32; the first future 15.625 to 33.225.
\item The index rose 95.9\%; the inverse lost 95.9\% on the day and 96.7\% since the start of the year.
\item An acceleration event on XIV: its intraday indicative value had fallen to 20\% or less of the previous close.
\item As a failure of hedge and leverage rebalancing in a concentrated, volatile market, like portfolio insurance.
\item \textbf{Named result.} Short futures at a quarter of capital earned 9.6\% a year with a 67\% drawdown; the inverse product grew to 16.5 times its
value and ended at zero when the first crash took the index from 17.3 to 63.6, facing a flow of 3.77 times its value at that close.
\item 4.8\% a year at 0.10 of capital, 9.6\% at 0.25; wiped out at 0.50 and 1.00.
\item Futures worth 3.77 times its value, twice the day's 189\%.
\item The index spends most of its time below its mean, so expected reversion lifts the futures.
\item By the largest one-day rise the book must survive, not by its roll-down.
\item Settlements, not index levels; 2018 and 2020 included; the roll and costs; several sizes.
\item ETP rebalancing anticipation.
\item Buyers of volatility futures pay it, for protection that pays when they need it.
\end{enumerate}
\end{solution}

\begin{solution}{iq:s2:the-volatility-index-complex:1}
The futures curve is usually in contango: each future converges down to the index as it nears expiry, and a long position of constant maturity keeps
selling the cheaper near contract and buying the dearer far one. Spikes are brief; the roll is constant.
\end{solution}

\begin{solution}{iq:s2:the-volatility-index-complex:2}
The market expects volatility to fall from a high level: a stress is under way. The roll now favours longs; short positions face the risk of a further
spike; calendar trades and small long positions are cheap to hold.
\end{solution}

\begin{solution}{iq:s2:the-volatility-index-complex:3}
Weight the first two contracts each close by the days left so that the maturity is constant, carry yesterday's weights through today's settlement, and
value an expiring contract at its final settlement. Errors come from expiry days, holidays, missing or zero settlements, and using index levels for
futures.
\end{solution}

\begin{solution}{iq:s2:the-volatility-index-complex:4}
By the largest one-day rise it must survive, with a margin; stress on a tripling of the index and on the products' flows at the close; limit size
so that the stress loss is a fraction of capital.
\end{solution}

\begin{solution}{iq:s2:the-volatility-index-complex:5}
At each close: the product's value from the index's return, the target exposure $L$ times the new value, the futures held after the move, and the
difference to trade; with checks on the settlement prices, the roll schedule and an acceleration trigger.
\end{solution}

\begin{solution}{iq:s2:the-volatility-index-complex:6}
As derived in the chapter: $LA(1 + Lr) - LA(1 + r) = A\,L(L - 1)\,r$. Over two days the product earns $(1 + Lr_1)(1 + Lr_2) - 1 = L(r_1 + r_2) +
L^2 r_1 r_2$, while $L$ times the index's two-day return is $L(r_1 + r_2) + L r_1 r_2$: they differ by $L(L - 1)r_1 r_2$.
\end{solution}
